% frontmatter/zusammenfassung.tex

\chapter*{Zusammenfassung}
\addcontentsline{toc}{chapter}{Zusammenfassung}

Diese Projektarbeit beschreibt das Fine-Tuning des \ac{vla}-Modells
\textbf{NVIDIA GR00T N1.6} auf einem humanoiden Roboter des Typs
\textbf{Unitree G1} mit \textbf{DEX3-Hand} für Block-Stacking-Aufgaben.
Mithilfe von \ac{il} auf einem aufgezeichneten Demonstrations-Datensatz wird das
vortrainierte Basismodell an die spezifische Aufgabe und Embodiment-Konfiguration
angepasst.

Die Trainingsinfrastruktur basiert auf einem Docker-Container mit CUDA~12.8 und
wird sowohl lokal als auch auf dem \ac{hpc}-Cluster KISSKI der GWDG Göttingen
(NVIDIA A100 / H100) ausgeführt.
Ziel ist eine möglichst hohe Erfolgsrate beim Block-Stacking unter realen oder
simulierten Bedingungen bei gleichzeitig effizienter Nutzung der verfügbaren
\ac{gpu}-Ressourcen.

\bigskip
\noindent\textbf{Schlüsselwörter:}
Humanoider Roboter,
\ac{vla},
Fine-Tuning,
\ac{il},
Unitree G1,
NVIDIA GR00T,
Block-Stacking

\vspace{2.5cm}

\chapter*{Abstract}
\addcontentsline{toc}{chapter}{Abstract}

This project report describes the fine-tuning of the \textbf{NVIDIA GR00T N1.6}
\ac{vla} model on a \textbf{Unitree G1} humanoid robot with \textbf{DEX3 hand}
for block-stacking tasks.
Using \ac{il} on a recorded demonstration dataset, the pre-trained foundation
model is adapted to the specific task and embodiment configuration.

The training infrastructure is based on a Docker container running CUDA~12.8,
executed both locally and on the KISSKI \ac{hpc} cluster of GWDG Göttingen
(NVIDIA A100 / H100).
The goal is a high success rate in block-stacking under real or simulated
conditions while making efficient use of available \ac{gpu} resources.

\bigskip
\noindent\textbf{Keywords:}
Humanoid Robot,
\ac{vla},
Fine-Tuning,
\ac{il},
Unitree G1,
NVIDIA GR00T,
Block-Stacking

\cleardoublepage
